\subsection{Existing Approaches}

Among the reviewed literature, three works are most directly relevant to this thesis, representing the closest prior work to the control problem being addressed. \\

The paper, Agile and Cooperative Aerial Manipulation of a Cable-Suspended Load, by TU Delft \cite{tudelft} proposes a centralised kinodynamic planning framework that this thesis takes as its primary performance benchmark. A receding-horizon optimal control problem is solved online at 10\,Hz over the full coupled load-cable-quadrotor dynamics, enforcing thrust limits, cable tautness, and collision avoidance directly as path constraints. Each quadrotor tracks its assigned trajectory using an onboard INDI controller that compensates for cable forces through IMU measurements, removing the need for any sensors on the load. Real-world experiments with three and four quadrotors achieve speeds exceeding 5\,m/s and accelerations exceeding 8\,m/s\textsuperscript{2}, at least eight times the agility of prior force-based methods, including flight through narrow passages. The primary limitation is that solve time grows exponentially with team size, and the framework assumes a fixed team throughout the mission with no support for dynamic agent changes. \\

Swarm-Based Dissipative Control for Cooperative Transport by Quan et al.\ \cite{quan2025swarm} experimentally demonstrates robustness to mid-mission agent dropout. The remaining UAVs self-organize and continue the trajectory with minimal payload swing without replanning, and further validates performance under varying drone capabilities, cable length uncertainty, and payload weight variation. This is achieved through a decentralized energy-dissipation controller in which each UAV independently damps relative motion between itself and the payload using only local measurements and no inter-agent communication. Any disturbance, including a UAV suddenly departing, injects excess kinetic energy into the system, which the dissipative law passively removes. The critical limitation is that this robustness does not extend to agent addition, as a joining drone receives no information about force sharing or attachment configuration, leaving dynamic team expansion unaddressed. \\

Tension-Based Reconfigurable Multi-Agent Formation by Bacheti et al.\ \cite{tensionbased} proposes a load-centric control framework in which the resultant cable tension acting on the payload is treated as a virtual control input. The controller computes the desired resultant tension through feedback linearisation, which is then distributed among the UAVs. Each agent tracks corresponding position and attitude references to realize the required cable forces, forming a hierarchical control structure linking load dynamics to individual vehicle actuation. A leader-follower strategy is introduced to enable formation reconfiguration, allowing the positions of the agents to be adjusted during operation, for example to navigate constrained environments such as narrow corridors. Experimental results demonstrate successful trajectory tracking and stable formation reconfiguration across multiple scenarios, including contraction and restoration of the formation. Although this approach is limited to agent reconfiguration, this idea could be beneficial when attempting addition or removal of agents. However, the system relies on simplified modeling assumptions, including a point-mass payload and rigid, always-taut cables, and does not incorporate adaptive force redistribution under uncertainty.
